% Verbatim copy of /Users/alingalatan/workspaces/grothe/handoff/research-degree5/ai-notes/chatgpt-thread2-remote_cubic_and_linear_shell.tex (file dated 2026-11-05). Only this line was added.
% output/r456/remote_cubic_and_linear_shell.tex
\documentclass[11pt]{article}
\usepackage[T1]{fontenc}
\usepackage{lmodern,microtype,amsmath,amssymb,mathtools,booktabs}
\usepackage[a4paper,top=23mm,bottom=23mm,left=25mm,right=25mm,headheight=14pt]{geometry}
\usepackage{xcolor,fancyhdr}
\usepackage[colorlinks=true,linkcolor=ink,urlcolor=ink]{hyperref}
\definecolor{ink}{HTML}{233B53}\definecolor{muted}{HTML}{505C65}
\pagestyle{fancy}\fancyhf{}\fancyhead[L]{\small\color{muted} Remote cubics and a linear shell}
\fancyhead[R]{\small\color{muted} 5 November 2026}\fancyfoot[C]{\small\thepage}
\renewcommand{\headrulewidth}{0.3pt}
\setlength{\parindent}{0pt}\setlength{\parskip}{4pt}
\setlength{\abovedisplayskip}{5pt plus 2pt minus 1pt}
\setlength{\belowdisplayskip}{5pt plus 2pt minus 1pt}
\newcommand{\E}{\mathbb E}\newcommand{\C}{\mathbb C}\newcommand{\R}{\mathbb R}
\newcommand{\J}{\mathcal J}\newcommand{\ind}{\mathbf 1}
\newcommand{\head}[1]{\par\smallskip{\large\bfseries\color{ink}#1}\par\smallskip}
\newcommand{\subhead}[1]{\par\smallskip{\bfseries\color{ink}#1}\enspace}
\newcommand{\qed}{\hfill\(\square\)}
\hypersetup{pdftitle={Remote cubics and a linear shell in Gaussian domination}}
\begin{document}
\thispagestyle{plain}
{\LARGE\bfseries\color{ink} Remote cubics and a linear shell\\[2pt] in Gaussian domination\par}
\vspace{3pt}{\color{muted} A fourth-order expansion and an all-amplitudes theorem for disjoint drivers\par}
\medskip
Let \(T=|z_1|^2\sim{\rm Exp}(1)\), \(a=\ell-2\beta>0\), \(f(t)=a+\beta t\),
\(r(t)=\sqrt t f(t)\), and \(L=z_1f(T)\). Write
\[
 F(G)=\E|L+G|+\E|L-G|,\qquad B(v)=2\E|L+\sqrt vW|,\qquad
 \Lambda(G)=F(G)-\|G\|_2^2/\varepsilon .
\]
The application box is \(.74\le\varepsilon\le.76,\ .76\le\ell\le.90,\ .12\le\beta\le.16\).
Set \(Q_t=q+(t-1)U\), where \(q\) is a remote \((2,1)\) Wick cubic and \(U\) is
a remote holomorphic linear form. Write \(v_q=\E|q|^2,\ u=\E|U|^2\).

\subhead{Main theorem.}
Throughout the box, if \(q\) and \(U\) are independent, then at every amplitude
\[
 \boxed{\quad
 \Lambda(q+(T-1)U)\le\Lambda(q)
       -\left(\varepsilon^{-1}-\frac{199}{200\ell}\right)u .
 \quad} \tag{1}
\]
The coefficient in parentheses is positive, since \(\ell\ge\varepsilon\). In particular
the given remote-only bound implies K for the entire disjoint case, including
\(U=z_2,\ q=z_3(|z_3|^2-2)\). The increment bound (1) holds for any independent
square-integrable remote offset \(q\). Pages 5--8 give the proof, including the
finite bounds for the whole parameter box.

\head{R4. The quadratic term}
Orthogonality of the first and third remote chaoses gives
\[
 m_2(t)=\E|Q_t|^2=v_q+u(t-1)^2,\qquad \|Q_T\|_2^2=v_q+u .
\]
For a comparator of variance \(b\lambda^2\), the coefficient of \(\lambda^2\) in
\(F(\lambda Q_T)-B(b\lambda^2)\) is
\[
 C_2(b)=\tfrac12\int_0^\infty \frac{e^{-t}}{r(t)}[m_2(t)-b]\,dt.
 \tag{2}
\]
At the matched variance \(b=v_q+u\), this becomes
\( \frac u2\int_0^\infty t(t-2)e^{-t}/r(t)\,dt\), strictly negative when \(u>0\):
\(t(t-2)\) has exponential mean zero and changes from negative to positive at 2,
while \(1/r(t)\) decreases strictly. The normalization in the question, \(v_q+2u=1\)
with comparator \(\lambda W\), subtracts the additional positive amount
\(\frac u2\int_0^\infty e^{-t}/r(t)\,dt\) from the matched coefficient.
Thus every fixed genuinely mixed direction satisfies strict small-amplitude
domination under both conventions, regardless of sharing between \(q\) and \(U\).
With total variance normalized to one, the conclusion is uniform on subfamilies
with \(u\ge\delta>0\): the integral giving the strictly negative quadratic factor
is continuous on the compact box and the \(O(\lambda^4)\) bound in the proof
below is uniform for bounded fourth moments. Fixed-order Gaussian hypercontractivity
supplies that moment bound independently of dimension.

\newpage
\head{The fourth coefficient and its regularization}
For any polynomial \(P\) of degree at most four define the convergent functional
\[
 \J[P]=\int_0^\infty
 \frac{e^{-t}P(t)/f(t)^3-P(0)/a^3}{t^{3/2}}\,dt .
\]
Let \(m_4(t)=\E|Q_t|^4\). The expansion is
\[
 F(\lambda Q_T)-B(b\lambda^2)
 =C_2(b)\lambda^2+\frac{\lambda^4}{32}\J[m_4-2b^2]+o(\lambda^4).
 \tag{3}
\]
With \(s=t-1\), its full input polynomial is
\[
 \begin{aligned}
 m_4(t)={}&\E|q|^4+4s\,\operatorname{Re}\E[|q|^2q\bar U]\\
 &+s^2\big(4\E[|q|^2|U|^2]+2\operatorname{Re}\E[q^2\bar U^2]\big)\\
 &+4s^3\operatorname{Re}\E[|U|^2q\bar U]+2u^2s^4 .
 \end{aligned}
\]
The odd powers need two further mixed moments. For instance, changing the sign
of a real \(c\) in the first shared example preserves all the listed even data.

\subhead{Proof of (3).}
Set \(H(z)=\E_\theta|1+ze^{i\theta}|\). Its series at zero is
\(1+z^2/4+z^4/64+O(z^6)\), by multiplying the two square-root series and
averaging the angle. For \(z\ge0\), the derivative of \(H(z)-1-z^2/4\) is
bounded in absolute value by a constant times \(z^3\). Use the series near zero and
the 1-Lipschitz property elsewhere. Thus after subtracting the quadratic term,
replacing \(|Q_t|\) by \(|Q_0|\) in the angular mean costs at most
\[
 C\lambda^4 r(t)^{-3}\,t |U|(|Q_t|+|Q_0|)^3 .
\]
This is integrable after division by \(\lambda^4\), including at \(t=0\).
Dominated convergence gives the ordinary quartic contribution
\(\frac1{64}\int e^{-t}(m_4(t)-m_4(0))/r(t)^3\,dt\) to one mean.

For the constant shift, \(g(x)=\E|L+x|\) is smooth and radial. The planar density of
\(L\) is \(h(|x|^2)/\pi\), where
\(h(r(t)^2)=e^{-t}/[f(t)(a+3\beta t)]\); it is smooth and rapidly decreasing.
Radial Taylor expansion gives coefficients \(\Delta g(0)/4\) and \(\Delta^2g(0)/64\).
The first Laplacian is \(\E|L|^{-1}\). Using \(\Delta|x|=|x|^{-1}\) and
integrating by parts gives
\[
 \Delta^2g(0)=4\int_0^\infty \frac{h'(y)+yh''(y)}{\sqrt y}\,dy
 =2\int_0^\infty\frac{h'(y)}{\sqrt y}\,dy=\J[1].
\]
For the last equality, integrate from \(y=r(\delta)^2\), change variables
\(y=r(t)^2\), and subtract \(2h(0)/r(\delta)=2/(a^3\sqrt\delta)+O(\sqrt\delta)\)
from the divergent integral of \(e^{-t}/r(t)^3\). The Taylor remainder divided by
\(|x|^4\) is bounded globally and tends to zero at zero, so averaging over \(Q_0\)
is justified by its fourth moment. Doubling the mean and subtracting the constant
Gaussian shift proves (3).\qed

\newpage
\head{Quartic signs for the three examples}
Write \(J_i=\J[(t-1)^i]/(\sqrt\pi/a^3)\) and use \(d=1\) for matched variance,
\(d=2\) for the normalization in the question. For the raw examples \(U=z_2\),
set \(b=v_q+d\). To impose \(v_q+2u=1\), divide the whole raw direction by
\(\sqrt{v_q+2}\). The corresponding quartic coefficient is the raw one divided by
\((v_q+2)^2\), preserving its sign.

Put \(R=|c|\), \(p=\operatorname{Re}c/R\) when \(R>0\). For the first example
\(q=cz_2(|z_2|^2-2)\) we have \(v_q=2R^2\). Its quartic sign is the sign of
\[
 \begin{aligned}
 P_1={}&2(J_4-d^2J_0)+8J_3Rp
  +[16(1+2p^2)J_2-8dJ_0]R^2\\
 &+128J_1R^3p+168J_0R^4 .
 \end{aligned}
\]
Indeed \(\E[T^2(T-2)^j]\), for \(j=0,\ldots,4\), equals \(2,2,8,32,176\).
For the second example \(q=cz_2(|z_3|^2-1)\), \(v_q=R^2\) and the corresponding
sign is the sign of
\[
 P_2=2(J_4-d^2J_0)+[(4+8p^2)J_2-4dJ_0]R^2
          +16J_1R^3p+16J_0R^4 .
\]
Here the first four centered moments of \(|z_3|^2\) are \(0,1,2,9\).

\textbf{Matched comparator.} Both signs are strictly negative for every complex
\(c\) throughout the box. The proof is on the next page.

\textbf{Comparator as written.} For real \(c\) the sign in each shared example is
positive between two roots, zero at the roots, and negative outside. The brackets
below hold throughout the box; the exact roots are those of the displayed polynomials.
\[
 \begin{array}{lcc}
 \toprule & \text{negative root} & \text{positive root}\\ \midrule
 q=cz_2(|z_2|^2-2)&(-1/5,-1/10)&(7/10,4/5)\\
 q=cz_2(|z_3|^2-1)&(-1/2,-2/5)&(9/10,1)\\ \bottomrule
 \end{array}
\]
For complex \(c\), the formulas give a direct phase test. At \(R>0\), write either
as \(A+Bp-Dp^2\), with \(D>0\). Its sign is that of
\(B^2+4DA-(2Dp-B)^2\). At \(c=0\) both signs are positive under this convention.
Thus the requested coefficient has no universal sign.

\textbf{Disjoint example.} For \(q=z_3(|z_3|^2-2)\), \(v_q=2\) and
\(m_4=176+8(t-1)^2+2(t-1)^4\). The coefficient has sign
\((176-2(2+d)^2)J_0+8J_2+2J_4<0\) for either convention.
More generally, for independent \(q,U\) the matched sign is the sign of
\[
 (\E|q|^4-2v_q^2)J_0+4v_qu(J_2-J_0)+2u^2(J_4-J_0),
\]
which is strictly negative if \(u>0\). A circular cubic has nonnegative excess
fourth moment: apply the real Wiener
\href{https://publications.uni.lu/bitstream/10993/13581/1/NouPecRei_AOP.pdf}{product formula}
to \(\operatorname{Re}q\)
and use circularity. For a nonzero cubic it is strictly positive: in the real
product formula the squared contraction of order two is nonzero.

\newpage
\head{Verification of the quartic signs}
If \(Y\sim\Gamma(3,1)\), \(z=\beta/a\), and \(k=1+zY\), then
\[
 \frac{\J[t^j]}{\sqrt\pi/a^3}
 =\frac{\Gamma(j-1/2)}{\sqrt\pi}\E k^{1/2-j},\qquad j=0,\ldots,4.
\]
This follows from \(f(t)^{-3}=a^{-3}\E e^{-ztY}\); for \(j=0\) integrate
\((e^{-kt}-1)t^{-3/2}\) to get \(-2\sqrt{\pi k}\).
Dividing by the positive factor \(\E\sqrt k\), the \(J_i\) are averages of the
following polynomials at \(h=1/k\in[0,1]\), with the law tilted by \(\sqrt k\):
\[
 \begin{aligned}
 &-2,\quad 2+h,\quad -2-2h+h^2/2,\quad 2+3h-3h^2/2+3h^3/4,\\
 &-2-4h+3h^2-3h^3+15h^4/8 .
 \end{aligned}
\]
In particular \(J_1,J_3>0,\ J_2,J_4<J_0<0\), and \(J_4-4J_0>0\).

For \(P_2\) at \(d=1\), maximizing the quadratic expression over all real \(p\)
bounds it by
\[
 2(J_4-J_0)+4(J_2-J_0)R^2+(16J_0-8J_1^2/J_2)R^4 .
\]
The last coefficient is negative: in the tilted law, put \(m=\E h,\ n=\E h^2\).
The sign of \(2J_0J_2-J_1^2\) is that of \(4+4m-2n-m^2>0\).
For \(P_1\) at \(d=1\), put \(D=-J_2\). Completing the square in \(p\) gives
\(DP_1\le A R^4+B R^2+C_0\), where
\[
 A=168J_0D+128J_1^2,\quad B=(16J_2-8J_0)D+16J_3J_1,\quad
 C_0=2(J_4-J_0)D+J_3^2/2 .
\]
The checker verifies \(A,C_0<0,\ 4AC_0>B^2\) for \(2/11\le z\le4/11\).

At \(d=2\), the coefficients in increasing order have signs \(+,{\ge}0,-,+,-\).
This gives one negative root.
For the positive side, let \(L\) be the left endpoint of the claimed positive
bracket. The checker verifies positive Bernstein coefficients on \([0,L]\),
strictly negative nonconstant coefficients after shifting by \(L\), and the asserted
endpoint signs. This proves existence, uniqueness and the sign classification.

\subhead{Exact enclosures.}
The script \texttt{verify\_core.mjs} integrates the five positive moments on
\([0,32]\) by composite Simpson with mesh \(1/128\). For the integrand
\(\tfrac12 y^2e^{-y}(1+zy)^{1/2-j}\), the fourth derivative is bounded by 2000:
all derivatives of both factors through order four have absolute value below 11
there, so Leibniz gives \(16\cdot11^2<2000\). The error is below \(1/750000\),
and the omitted gamma tail is below \(10^{-8}\) by comparison with
\(\frac12 y^2(1+y)e^{-y}\). Monotonicity of each moment in \(z\) encloses full
profile intervals using their endpoints. The replay covered 40 intervals.
Every operation used for the certificate rounds outwards on a \(2^{-40}\) grid.

\newpage
\head{R5. Why the pointwise test fails, and what replaces it}
For any circular \(Y\) with a fourth moment, the angular expansion on page 2 gives
\[
 \E|r+Y|=r+\frac{\E|Y|^2}{4r}+\frac{\E|Y|^4}{64r^3}+o(r^{-3})
 \qquad(r\to\infty).
\]
Indeed \(|H(z)-1-z^2/4|\le C z^4\) globally, so dominated convergence applies.
Adding an independent proper Gaussian \(X\) preserves the excess fourth moment
\(\E|Y|^4-2(\E|Y|^2)^2\). Positive excess therefore reverses the proposed
comparison for all sufficiently large \(r\).

This occurs even for Gaussian scale mixtures with fourth moment arbitrarily close
to the Gaussian value. For instance take
\(q_N=N^{-1/2}\sum_{j=1}^N X_j(|Y_j|^2-1)\) from independent complex Gaussians.
Conditional on the \(Y_j\) it is proper Gaussian with variance
\(V_N=N^{-1}\sum (|Y_j|^2-1)^2\), of mean one and variance \(8/N\).
Thus \(\E|q_N|^4=2+16/N\). This disproves the pointwise claim for every proposed
fourth-moment cap \(\kappa>2\).

\subhead{Proof of the averaged theorem (1).}
Let \(M_w(x)=\E|x+\sqrt w W|\) and \(K_w(x)=\E|x+\sqrt w W|^{-1}\).
The planar heat equation gives \(\partial_w M_w=K_w/4\). For a deterministic offset
\(y\in\C\), integrating in the variance of \(U\) gives the increase in a single mean
\[
 \frac14\int_0^u \mathcal I_s(y)\,ds,\qquad
 \mathcal I_s(y)=
 \int_0^\infty e^{-t}(t-1)^2
       \E_\theta K_{s(t-1)^2}\bigl(r(t)e^{i\theta}+y\bigr)\,dt .
 \tag{4}
\]
One may first start at positive variance and pass down to zero, using positivity
and continuity of \(M\). Independence permits averaging (4) over \(y=q\) and
using it for both signs. Orthogonality adds exactly \(u/\varepsilon\) to the cost.
It thus suffices to bound \(\mathcal I_s(y)<1.99/\ell\), uniformly in \(s,y\).
Pages 6--8 prove this.

\subhead{A shorter estimate at the three application rows.}
The planar density of the weighted visible measure, up to \(1/\pi\), at radius
\(R(x)=x(a+\beta x^2)\) is
\[
 h(x)=\frac{(x^2-1)^2e^{-x^2}}{(a+\beta x^2)(a+3\beta x^2)} .
\]
Using \(K_w(b)\le 6/(5b)\), majorize \(h\) by
\(\widehat h(x)=\sup_{y\ge x} h(y)\). Convolution of two nonnegative radial decreasing
functions takes its maximum at zero, by the disk layer-cake representation.
Consequently \(\mathcal I_s(y)\le(6/5)\,2\int_0^\infty\widehat h(x)R'(x)\,dx\).
The last integral including the factor 2 is \(<2.14\) at each actual row.
The checker uses the right and left endpoints of mesh-\(1/128\) cells to enclose
\((x^2-1)^2\) and \(e^{-x^2}/[(a+\beta x^2)(a+3\beta x^2)]\), respectively, then
backward maxima. Beyond \(x=4\), \(h\) decreases and
\((2/a)\int_4^\infty x^4e^{-x^2}\,dx<1/1000\).
This already gives \(F(Q_T)-F(q)<(321/250)u\) when \(u>0\), sufficient for (strict)
penalized domination at the three rows.

\newpage
\head{Three bounds for the inverse-radius kernel}
For \(r,w>0\), \(x=r^2/(2w)\), Gaussian integration gives
\(K_w(r)=\sqrt{\pi/w}\,e^{-x}I_0(x)\).
We use
\[
 r K_w(r)\ \le\ \min\left(\frac{237}{200},\
           \frac{\sqrt\pi\,r}{\sqrt{w+r^2}},\
           1+\frac{w}{2r^2}\right).
 \tag{5}
\]
Here are details to make the numerical certificate on the next pages reproducible.
Put \(p=I_1/I_0\), \(m=2x(1-p)\),
\(U(x)=\sqrt{2\pi x}e^{-x}I_0(x)\). The Bessel recurrences give
\[
 p'=1-p/x-p^2,\qquad
 m'=2-2m+m^2/(2x),\qquad U'/U=(1-m)/(2x).
\]
These identities also follow by the series for \(I_0\);
see \href{https://dlmf.nist.gov/10.29}{NIST DLMF \S10.29} and
\href{https://dlmf.nist.gov/10.32}{\S10.32}.

At \(x=1\), \(I_1(1)<I_0(1)/2\) by the series.
Since \(m'>0\) at every crossing of \(1\), \(U\) decreases for \(x\ge1\).
On \(0<x\le1\), use \(I_0(x)\le e^{x^2/4}\); the bound
\(\sqrt{2\pi x}\exp(-x+x^2/4)\) increases there. Its value at 1 is
\(\sqrt{2\pi}e^{-3/4}<237/200\).

For the middle bound, \(p< b(x)=2x/(1+2x)\). This holds near zero by the series,
and at a crossing \(p'-b'=-3/(1+2x)^2<0\), preventing an upward crossing.
Thus \(\sqrt{1+2x}\,e^{-x}I_0(x)\) is decreasing from 1, as claimed.

For the last bound put \(v(x)=1+1/(4x)\). The constant bound already gives
\(U(x)\le v(x)\) for \(0<x\le4/3\). On the two intervals
\([4/3,7/4]\) and \([7/4,9/4]\) use monotonicity of \(U\) and the endpoint
checks \(U(4/3)<8/7,\ U(7/4)<10/9\); these are enclosed by the series in the checker.
For \(x\ge9/4\) write
\(\mathcal L y=y''+2y'+y/(4x^2)\). Then \(\mathcal L U=0\) and
\(\mathcal L v=(9-4x)/(16x^3)\le0\).
Consequently
\[
 \big[e^{2x}(vU'-Uv')\big]'=-e^{2x}U\mathcal L v\ge0 .
\]
The ratio \(U/v\) has its supremum at an endpoint of this half-line. At \(9/4\)
it is below one, and at infinity it tends to one. The latter limit follows also
by scaling the angle in \(e^{-x}I_0(x)=(2/\pi)\int_0^{\pi/2}e^{-2x\sin^2 y}\,dy\)
and using \(\sin y\ge2y/\pi\). This proves (5).\qed

\newpage
\head{The full-box potential bound}
Divide all lengths by \(\ell\). Write \(p=\beta/\ell\in[2/15,4/19]\) and
\(d=\sqrt{s}/\ell\). The normalized radius, its derivative, and the density
of the weighted visible measure, up to \(1/\pi\), are
\[
 R_p(x)=x[1+p(x^2-2)],\quad R_p'(x)=1+p(3x^2-2),\quad
 g_p(x)=\frac{(x^2-1)^2e^{-x^2}}{[1+p(x^2-2)][1+p(3x^2-2)]}.
\]
The required bound on the normalized potential is \(1.99\), and the original
potential equals \(1/\ell\) times the normalized one.

\subhead{Disks with variable Gaussian variance.}
For \(0\le x\le X\) the Gaussian standard deviation \(d|x^2-1|\)
lies between \(z=d_0(1-X^2)_+\) and \(d_1\max(1,|X^2-1|)\) if
\(d_0\le d\le d_1\). Take the supremum of \(K_w\) over all variances in this
interval. It is radially decreasing because every \(K_w\) is. Its integral on
any translated disk of radius \(R\) is at most its integral on the disk centered
at zero. This follows by integrating areas of intersections of disks.

Divide that integral by \(2\pi\) and put \(C=237/200\). An upper bound is
\[
 D(R,X;d_0,d_1)
 =C(R-m)+1.773\big(\sqrt{z^2+m^2}-z\big)
 -(C-1)(R-F)_+^2/R , \tag{6}
\]
where \(m=\min(R,17z/20)\), \(F=(5/3)d_1\max(1,|X^2-1|)\).
For \(d_1=\infty\) the last term is zero. To see (6), integrate \(rK_w(r)\)
in the radial coordinate. On \([0,m]\), use the middle bound of (5) with
\(\sqrt\pi<1.773\). On the remaining interval first use \(C\). Beyond \(F\),
the last bound of (5) is at most \(1+(C-1)F^2/r^2\), since
\((C-1)(5/3)^2>1/2\). Integrating the gain over \([F,R]\) gives
\((C-1)(R-F)^2/R\) when \(F<R\). These two improved intervals do not overlap.

\subhead{Positive disk decomposition.}
On \(0\le x\le4\) take cells \([l_j,x_j]\) of length \(1/1024\). For each compact
profile interval, bound \(g_p\) above on each cell: use the larger endpoint value
of \((x^2-1)^2\), the left endpoint for the decreasing exponential, and the lower
endpoint bounds for both positive denominator factors. Take backward maxima \(H_j\)
of the resulting heights and put \(H_{n+1}=0\). Then the contribution of this
range is bounded by
\[
 2\sum_j (H_j-H_{j+1})D(R_j,x_j;d_0,d_1), \tag{7}
\]
where \(R_j\) bounds \(R_p(x_j)\) from above. Indeed the step majorant is a positive
sum of disk indicators. On each actual disk apply the variance interval above and
then extend the disk to radius \(R_j\). This last step uses the positive integral
before taking its explicit upper bound (6).

\newpage
\head{Tail and exact verification for (1)}
On \(x\ge4\), \(g_p(x)\) decreases. By the constant bound in (5), a radially
decreasing majorant for this tail is obtained by extending the value \(g_p(4)\)
as a constant over the inner disk. The normalized tail potential is at most
\[
 2C\left[g_p(4)R_p(4)+
   \int_4^\infty\frac{(x^2-1)^2e^{-x^2}}{1+p(x^2-2)}\,dx\right]<\frac1{5000}.
\]
For example \(R_p(4)<16\), the denominator product in \(g_p(4)\) is larger than
20, \(e^{-16}<1.3\cdot10^{-7}\), and
\(\int_4^\infty x^4e^{-x^2}\,dx\le(290/8)e^{-16}\), by \(x^4\le x^5/4\).

The companion \texttt{verify\_core.mjs} bounds (7) plus this tail using integers.
For the compact range split \(p\) and \(d\) into intervals, enlarging the outer
profile endpoints by directed rounding. Evaluate (6) using a lower bound on \(z\),
an upper bound on \(F\), and an upper bound on the square root. The small integration
cutoff \(m\) may round down. For \(d\ge1\), use (6) with \(d_0=1,d_1=\infty\).
The resulting strict upper bound, in units of \(2^{-40}\), is
\[
 \frac{2188027284403}{1099511627776}<1.99 .
\]
The replay covered 12 profile intervals and 278 variance intervals, including
the unbounded tails. This proves (1) via (4). For \(e^{-x}\) upper bounds the
code takes reciprocals of lower partial sums of \(e^x\). Every integer square
root uses exact comparisons after its initial guess. Run
\texttt{node verify\_core.mjs}; no external package is required.\qed

\subhead{What is used from the remote problem.}
The proof above gives (1) directly, regardless of the law of the independent
\(q\). Applying the remote-only \(\Lambda(q)\le S\) supplied in the question
finishes the disjoint case. A remote-only estimate \(S+\eta\) would transfer
with the same error and the displayed strict improvement from the shell.
Independence enters when the variable \(y\) in (4) can be held fixed during
Gaussian integration in the driver \(U\). For arbitrary shared \(q,U\), the
small-amplitude theorem on pages 1--2 applies. All-amplitudes K and the universal
sign of the \emph{matched} quartic for arbitrary shared tensors remain open here.

For a disjoint Gaussian scale mixture \(q\stackrel d=\sqrt V W'\), the stronger
comparison \(F(Q_T)\le B(v_q+u)\) also follows from the previously checked
conditional Gaussian theorem. Given \(V,z_1\), the variance is \(V+u(T-1)^2\).
Apply that theorem and then Jensen to the concave \(B\).
For arbitrary disjoint cubics, (1) settles the penalized comparison; a same-norm
comparison at every amplitude would be stronger.

\newpage
\head{R6. What the second-moment Laplace method can prove}
Suppose only that \(\E[|G|^2\mid T]=vT+u(T-1)^2\). The signed Jensen step of R1
gives \(F(G)/2\le\E\sqrt{r(T)^2+vT+u(T-1)^2}\).
For every \(u>0\), the variable under the square root is bounded below by a
strictly positive constant: its continuous expression cannot vanish at any \(t\)
and tends to infinity. Its Laplace transform thus decays exponentially. With
\(w=v+u\), the Gaussian transform instead satisfies
\[
 \E e^{-\xi|L+\sqrt wW|^2}
 =\frac{\E e^{-\xi r(T)^2/(1+\xi w)}}{1+\xi w}
 \ \sim\ \frac{\E e^{-r(T)^2/w}}{\xi w}.
\]
Therefore the R1 Laplace order fails for all sufficiently large \(\xi\). A change
of variables cannot prove that order.

Even the mean comparison can fail under the conditional moment hypothesis alone.
Take \(G=i(z_1/|z_1|)(T-1)\sqrt u\), assigning any value at \(z_1=0\).
This square-integrable charge-one function attains the Jensen envelope. The
coefficient of \(u\) in \(F(G)-B(u)\) at zero is
\[
 \int_0^\infty\frac{e^{-t}}{r(t)}[(t-1)^2-1/2]\,dt
 \ \ge\ \frac16\int_0^\infty\frac{e^{-t}}{r(t)}\,dt>0 .
\]
Indeed the probability density proportional to \(e^{-t}/r(t)\) is a positive
mixture of gamma densities of shape \(1/2\) and rates \(k\ge1\), by the Laplace
representation of \(1/(a+\beta t)\). In each component
\(\E(T-1)^2=1-1/k+3/(4k^2)\ge2/3\). The example is outside the cubic chaos.
Strict failure persists for sufficiently small positive \(v/u\) by continuity
if the perpendicular amplitude is changed to \(\sqrt{vT+u(T-1)^2}\).

\subhead{A positive region.}
Put \(R(w)=\E\sqrt T\sqrt{f(T)^2+w}\) and
\[
 d(w)=\frac12\int_0^\infty
 \frac{e^{-t}(t^2-3t+1)}{\sqrt t\sqrt{f(t)^2+w}}\,dt .
\]
The square-root tangent bound gives \(F(G)/2\le R(w)+u d(w)\). Thus the explicit
test \(u d(w)\le B(w)/2-R(w)\) suffices. In particular, uniformly on the box,
\[
 0<w,\qquad 0\le u\le\frac{w}{10000(1+w)^3}
 \quad\Longrightarrow\quad F(G)\le B(w).
\]
For the constants, the R1 Laplace comparison with \(D=1+\xi w\) replaces
\(f(y)^2\) by \(f(y/D)^2\). Restrict to \(1\le y\le2,\ 1/2\le\xi\le1\) and use
\(e^x-1\ge x\). Integration against \(\xi^{-3/2}\,d\xi/(2\sqrt\pi)\) gives
\[
 B(w)/2-R(w)\ge
 \frac{\sqrt2-1}{4\sqrt\pi}\,e^{-2-2\ell^2}
 \left(\frac{14}{3}a\beta+\frac{15}{4}\beta^2\right)\frac{w}{(1+w)^3}
 >\frac{w}{2400(1+w)^3}.
\]
Here \(1-1/D^2\ge1-1/D=\xi w/D\), producing the two moments of \(y\) in the
displayed coefficient. Also \(d(w)\le7\sqrt\pi/(8a)<4\). This proves the stated
range. The structural Gaussian and disjoint cases have the stronger results above.
\end{document}
